\begin{lemma}[A telescoping square-root sum bound and a logarithmic square bound]
The following two elementary inequalities hold.
\begin{enumerate}
\item For every real $a>0$ and every integer $M\ge 0$,
\begin{equation}\label{eq:sslsb-harm}
  \sum_{m=0}^{M-1}\frac{1}{\sqrt{a(m+1)}}\le\frac{2\sqrt{aM}}{a}.
\end{equation}
\item For every real $s$ with $0<s\le 8$,
\begin{equation}\label{eq:sslsb-log}
  |\log s|^2\le\frac{16}{\sqrt{s}}+(\log 8)^2 .
\end{equation}
\end{enumerate}
\end{lemma}
\begin{proof}
\emph{Part 1: the telescoping comparison.}
Fix $a>0$ and $M\ge 0$. We first prove the termwise estimate
\begin{equation}\label{eq:sslsb-term}
  \frac{1}{\sqrt{a(m+1)}}\le\frac{2\bigl(\sqrt{a(m+1)}-\sqrt{am}\bigr)}{a}
  \qquad(m\ge 0\text{ an integer}).
\end{equation}
Write $u=\sqrt{a(m+1)}$ and $w=\sqrt{am}$. Since $a>0$ and $m+1\ge 1>0$ we have
$a(m+1)>0$, so $u>0$ and $u^2=a(m+1)$; since $m\ge 0$ we have $am\ge 0$, so
$w\ge 0$ and $w^2=am$. Consequently
\begin{equation}\label{eq:sslsb-diff}
  u^2-w^2=a(m+1)-am=a ,
\end{equation}
and $w\le u$ because $am\le a(m+1)$ and the square root is monotone on
$[0,\infty)$. Both $u>0$ and $a>0$, so \eqref{eq:sslsb-term} is equivalent, after
multiplying both sides by the positive number $ua$, to
\[
  a\le 2u(u-w),
\]
and by \eqref{eq:sslsb-diff} this is in turn equivalent to
$u^2-w^2\le 2u^2-2uw$, i.e.\ to $0\le u^2-2uw+w^2=(u-w)^2$, which holds. This
proves \eqref{eq:sslsb-term}.

Next we evaluate the sum of the right-hand sides of \eqref{eq:sslsb-term}. Put
$g(k)=\sqrt{ak}$ for integers $k\ge 0$. For each integer $m\ge 0$,
\[
  \frac{2\bigl(\sqrt{a(m+1)}-\sqrt{am}\bigr)}{a}
   =\frac{2}{a}\bigl(g(m+1)-g(m)\bigr),
\]
so by the telescoping identity
$\sum_{m=0}^{M-1}\bigl(g(m+1)-g(m)\bigr)=g(M)-g(0)$ and by $g(0)=\sqrt{0}=0$,
\begin{equation}\label{eq:sslsb-tele}
  \sum_{m=0}^{M-1}\frac{2\bigl(\sqrt{a(m+1)}-\sqrt{am}\bigr)}{a}
   =\frac{2}{a}\bigl(g(M)-g(0)\bigr)=\frac{2\sqrt{aM}}{a}.
\end{equation}
Summing \eqref{eq:sslsb-term} over $m=0,\dots,M-1$ and using monotonicity of
finite sums together with \eqref{eq:sslsb-tele} gives
\[
  \sum_{m=0}^{M-1}\frac{1}{\sqrt{a(m+1)}}
   \le\sum_{m=0}^{M-1}\frac{2\bigl(\sqrt{a(m+1)}-\sqrt{am}\bigr)}{a}
   =\frac{2\sqrt{aM}}{a},
\]
which is \eqref{eq:sslsb-harm}.

\emph{Part 2: the logarithmic square bound.}
Fix $s$ with $0<s\le 8$. Since $8\ge 1$ we have $\log 8\ge 0$. We distinguish
two cases.

\emph{Case $s\le 1$.} Set $u=\sqrt{\sqrt{s}}$. As $s>0$ we have $\sqrt{s}>0$ and
hence $u>0$, and $u^2=\sqrt{s}$. From $\log\sqrt{x}=\tfrac12\log x$ for $x\ge 0$,
applied twice,
\begin{equation}\label{eq:sslsb-logu}
  \log u=\tfrac12\log\sqrt{s}=\tfrac14\log s .
\end{equation}
The inequality $\log x\le x-1$, valid for all $x>0$, applied to $x=u^{-1}>0$
gives $\log(u^{-1})\le u^{-1}-1$, i.e.\ $-\log u\le \frac{1}{u}-1\le\frac1u$.
Combining this with \eqref{eq:sslsb-logu} and multiplying by $4>0$,
\begin{equation}\label{eq:sslsb-logs}
  -\log s\le\frac{4}{u}.
\end{equation}
Because $0<s\le 1$ we have $\log s\le 0$, so $|\log s|=-\log s$ and
$-\log s\ge 0$. Both sides of \eqref{eq:sslsb-logs} are therefore nonnegative,
and squaring a valid inequality between nonnegative reals preserves it:
\[
  (-\log s)^2\le\Bigl(\frac{4}{u}\Bigr)^2=\frac{16}{u^2}=\frac{16}{\sqrt{s}},
\]
the last step by $u^2=\sqrt s$. Since $(\log 8)^2\ge 0$, we conclude
\[
  |\log s|^2=(-\log s)^2\le\frac{16}{\sqrt{s}}\le\frac{16}{\sqrt{s}}+(\log 8)^2 ,
\]
which is \eqref{eq:sslsb-log} in this case.

\emph{Case $s>1$.} Then $\log s\ge 0$, so $|\log s|=\log s$, and $s\le 8$ with
$s>0$ gives $\log s\le\log 8$ by monotonicity of the logarithm on $(0,\infty)$.
Both $\log s$ and $\log 8$ are nonnegative with $\log s\le\log 8$, so
$(\log s)^2\le(\log 8)^2$. Finally $16/\sqrt{s}\ge 0$ because $s>0$, whence
\[
  |\log s|^2=(\log s)^2\le(\log 8)^2\le\frac{16}{\sqrt{s}}+(\log 8)^2 .
\]
This is \eqref{eq:sslsb-log} in the remaining case, and the two cases exhaust
$0<s\le 8$.
\end{proof}
